% GENERATED FILE. Edit canonical lessons, then run npm run generate:books.
% canonical-source-hash: fnv1a64:90af29fd7f22a740
% canonical-lessons: RU-C488-bessonnitsa, RU-C488-khrapet, RU-C488-zevat, RU-C488-koshmar, RU-C488-dremat

\chapter{A Sleepless Night}
\label{ch:ru-a-sleepless-night}
\hlchaptermodality{488}

\begin{chapteropening}
\textbf{By the end of this chapter:} \emph{I can describe a bad night: insomnia, snoring and nightmares.}

Complete the last lesson of chapter 488: I can describe a bad night: insomnia, snoring and nightmares.
\end{chapteropening}

\section[\texorpdfstring{\ru{бессонница} (bessónnitsa)}{бессонница (bessónnitsa)}]{\ru{бессонница} (bessónnitsa) — insomnia}

\label{lesson:RU-C488-bessonnitsa}



\begin{quote}
Before the new one: say the Russian for a vaccination, then the Russian for to get infected.
\end{quote}

\subsection*{You'll want to know: \ru{бессонница}}
\textbf{\ru{бессонница}} (\emph{bessónnitsa}) — \textquotedblleft{}insomnia\textquotedblright{}.

\begin{quote}
\ru{Из}-\ru{за} \ru{стресса} \ru{у} \ru{меня} \ru{бессонница}.
\end{quote}

\begin{quote}
\emph{Iz-za stréssa u menyá bessónnitsa.}
\end{quote}

\begin{quote}
I've got insomnia because of stress.
\end{quote}

Feminine.

\begin{cousinweb}
\emph{\ru{бессонница}} is built on \emph{\ru{бессонный}} 'sleepless', from \emph{\ru{без}} 'without' and \emph{\ru{сон}} 'sleep', with \emph{-\ru{ница}}; EtymDB lists \emph{\ru{бессонный}} without a source and has no entry for \emph{\ru{бессонница}}.
\end{cousinweb}

\begin{grammarlens}[title={What you've built}]
One more word for telling, explaining and describing at B1.
\end{grammarlens}

\subsection*{Your turn}
\begin{itemize}
\raggedright
  \item \emph{Say it:} \emph{bessónnitsa}
  \item \emph{Say it:} \emph{bessónnitsa}, once more
  \item \emph{Say it:} a sentence of your own with \emph{bessónnitsa}
  \item \emph{Recall:} say the Russian for a vaccination, then the Russian for to get infected, then say \emph{bessónnitsa} again
\end{itemize}

\begin{tcolorbox}[breakable,colback=teal!4,colframe=teal!35!black,title={Before you move on}]
What does \ru{бессонница} mean? (\textquotedblleft{}Insomnia\textquotedblright{}.) Say it once more.
\end{tcolorbox}
\section[\texorpdfstring{\ru{храпеть} (khrapét)}{храпеть (khrapét)}]{\ru{храпеть} (khrapét) — to snore}

\label{lesson:RU-C488-khrapet}



\begin{quote}
Before the new one: say the Russian for to get infected, then the Russian for insomnia.
\end{quote}

\subsection*{You'll want to know: \ru{храпеть}}
\textbf{\ru{храпеть}} (\emph{khrapét}) — \textquotedblleft{}to snore\textquotedblright{}.

\begin{quote}
\ru{Мой} \ru{муж} \ru{так} \ru{храпит}, \ru{что} \ru{я} \ru{не} \ru{могу} \ru{заснуть}.
\end{quote}

\begin{quote}
\emph{Moy muzh tak khrapít, chto ya ne mogú zasnút.}
\end{quote}

\begin{quote}
My husband snores so loudly that I can't get to sleep.
\end{quote}

Imperfective; present \ru{храплю}, \ru{храпишь} (\ru{п} becomes \ru{пл}). The noun is \emph{\ru{храп}} — snoring.

\begin{cousinweb}
EtymDB lists \emph{\ru{храпеть}}, glossed 'to snore', without a source; it is usually taken as imitative.
\end{cousinweb}

\begin{grammarlens}[title={What you've built}]
One more word for telling, explaining and describing at B1.
\end{grammarlens}

\subsection*{Your turn}
\begin{itemize}
\raggedright
  \item \emph{Say it:} \emph{khrapét}
  \item \emph{Say it:} \emph{khrapét}, once more
  \item \emph{Say it:} a sentence of your own with \emph{khrapét}
  \item \emph{Recall:} say the Russian for to get infected, then the Russian for insomnia, then say \emph{khrapét} again
\end{itemize}

\begin{tcolorbox}[breakable,colback=teal!4,colframe=teal!35!black,title={Before you move on}]
What does \ru{храпеть} mean? (\textquotedblleft{}To snore\textquotedblright{}.) Say it once more.
\end{tcolorbox}
\section[\texorpdfstring{\ru{зевать} (zevát)}{зевать (zevát)}]{\ru{зевать} (zevát) — to yawn}

\label{lesson:RU-C488-zevat}



\begin{quote}
Before the new one: say the Russian for insomnia, then the Russian for to snore.
\end{quote}

\subsection*{You'll want to know: \ru{зевать}}
\textbf{\ru{зевать}} (\emph{zevát}) — \textquotedblleft{}to yawn\textquotedblright{}.

\begin{quote}
\ru{На} \ru{скучном} \ru{уроке} \ru{все} \ru{зевали}.
\end{quote}

\begin{quote}
\emph{Na skúchnom uróke vse zeváli.}
\end{quote}

\begin{quote}
Everyone was yawning in the boring lesson.
\end{quote}

Imperfective; perfective \ru{зевнуть} (once).

\begin{cousinweb}
\emph{\ru{зевать}} is inherited from Proto-Slavic \emph{zěvati} (EtymDB).
\end{cousinweb}

\begin{grammarlens}[title={What you've built}]
One more word for telling, explaining and describing at B1.
\end{grammarlens}

\subsection*{Your turn}
\begin{itemize}
\raggedright
  \item \emph{Say it:} \emph{zevát}
  \item \emph{Say it:} \emph{zevát}, once more
  \item \emph{Say it:} a sentence of your own with \emph{zevát}
  \item \emph{Recall:} say the Russian for insomnia, then the Russian for to snore, then say \emph{zevát} again
\end{itemize}

\begin{tcolorbox}[breakable,colback=teal!4,colframe=teal!35!black,title={Before you move on}]
What does \ru{зевать} mean? (\textquotedblleft{}To yawn\textquotedblright{}.) Say it once more.
\end{tcolorbox}
\section[\texorpdfstring{\ru{кошмар} (koshmár)}{кошмар (koshmár)}]{\ru{кошмар} (koshmár) — a nightmare}

\label{lesson:RU-C488-koshmar}



\begin{quote}
Before the new one: say the Russian for to snore, then the Russian for to yawn.
\end{quote}

\subsection*{You'll want to know: \ru{кошмар}}
\textbf{\ru{кошмар}} (\emph{koshmár}) — \textquotedblleft{}a nightmare\textquotedblright{}.

\begin{quote}
\ru{Ночью} \ru{мне} \ru{приснился} \ru{кошмар}.
\end{quote}

\begin{quote}
\emph{Nóchyu mne prisnílsya koshmár.}
\end{quote}

\begin{quote}
I had a nightmare in the night.
\end{quote}

Masculine. Colloquially also an exclamation: \emph{\ru{Какой} \ru{кошмар}!} — How awful!

\begin{cousinweb}
EtymDB lists \emph{\ru{кошмар}} without a source; it is usually traced to French \emph{cauchemar} 'nightmare'.
\end{cousinweb}

\begin{grammarlens}[title={What you've built}]
One more word for telling, explaining and describing at B1.
\end{grammarlens}

\subsection*{Your turn}
\begin{itemize}
\raggedright
  \item \emph{Say it:} \emph{koshmár}
  \item \emph{Say it:} \emph{koshmár}, once more
  \item \emph{Say it:} a sentence of your own with \emph{koshmár}
  \item \emph{Recall:} say the Russian for to snore, then the Russian for to yawn, then say \emph{koshmár} again
\end{itemize}

\begin{tcolorbox}[breakable,colback=teal!4,colframe=teal!35!black,title={Before you move on}]
What does \ru{кошмар} mean? (\textquotedblleft{}A nightmare\textquotedblright{}.) Say it once more.
\end{tcolorbox}
\section[\texorpdfstring{\ru{дремать} (dremát)}{дремать (dremát)}]{\ru{дремать} (dremát) — to doze}

\label{lesson:RU-C488-dremat}



\begin{quote}
Before the new one: say the Russian for to yawn, then the Russian for a nightmare.
\end{quote}

\subsection*{You'll want to know: \ru{дремать}}
\textbf{\ru{дремать}} (\emph{dremát}) — \textquotedblleft{}to doze\textquotedblright{}.

\begin{quote}
\ru{Дедушка} \ru{дремлет} \ru{в} \ru{кресле}.
\end{quote}

\begin{quote}
\emph{Dédushka drémlet v krésle.}
\end{quote}

\begin{quote}
Grandad is dozing in the armchair.
\end{quote}

Imperfective; present \ru{дремлю}, \ru{дремлешь} (\ru{м} becomes \ru{мл}). Perfective \ru{задремать} — to doze off.

\begin{cousinweb}
\emph{\ru{дремать}} is inherited from Proto-Slavic \emph{drěmati} (EtymDB).
\end{cousinweb}

\begin{grammarlens}[title={What you've built}]
One more word for telling, explaining and describing at B1.
\end{grammarlens}

\subsection*{Your turn}
\begin{itemize}
\raggedright
  \item \emph{Say it:} \emph{dremát}
  \item \emph{Say it:} \emph{dremát}, once more
  \item \emph{Say it:} a sentence of your own with \emph{dremát}
  \item \emph{Recall:} say the Russian for to yawn, then the Russian for a nightmare, then say \emph{dremát} again
\end{itemize}

\begin{tcolorbox}[breakable,colback=teal!4,colframe=teal!35!black,title={Before you move on}]
What does \ru{дремать} mean? (\textquotedblleft{}To doze\textquotedblright{}.) Say it once more.
\end{tcolorbox}
